\ifdefined\gkver\else\def\gkver{student}\fi
\documentclass[\gkver]{gaokaozhenti}
\begin{document}

\chapter{1995年全国}

\begin{enumerate}
\item
%% number: 1
%% paperName: 1995年全国高考第 1 题 3 分
%% typeId: 1
%% type: 单选题
%% chapter: 力物体的平衡
%% point: 物体系平衡
%% method: 无
%% score: 3
%% degree: 850
%% duplicateId: 0
%% body:
如图 \ref{1995全国01} 所示，两个物体 A 和 B，质量分别为 $M$ 和 $m$，用跨过定滑轮的轻绳相连，A 静止于水平地面上，不计摩擦，则 A 对绳的作用力与地面对 A 的作用力的大小分别为\xzanswer{A}
\onepicture[num=1, showlabel=false, label=1995全国01, w=4cm]{\includesvg{figs/01}}
\fourchoices[answer=A]
{$mg$，$(M-m)g$}
{$mg$，$Mg$}
{$(M-m)g$，$Mg$}
{$(M+m)g$，$(M-m)g$}

%% answer: A
%% memo:
\memoanswer{
由 B 受力平衡知，绳对 B 的拉力大小为 $T = mg$，所以绳对 A 的拉力大小为 $mg$，设地面对 A 的作用力大小为 $F$，由 A 受力平衡得，$F + mg = Mg$，解得 $F = (M - m)g$，故 A 正确。
}

\item
%% number: 2
%% paperName: 1995年全国高考第 2 题 3 分
%% typeId: 1
%% type: 单选题
%% chapter: 光学
%% point: 光的衍射
%% method: 无
%% score: 3
%% degree: 850
%% duplicateId: 0
%% body:
在观察光的衍射现象的实验中，通过紧靠眼睛的卡尺测脚形成的狭缝，观看远处的日光灯管或线状白炽灯丝（灯管或灯丝都要平行于狭缝），可以看到\xzanswer{C}
\fourchoices[answer=C]
{黑白相间的直条纹}
{黑白相间的弧形条纹}
{彩色的直条纹}
{彩色的弧形条纹}

%% answer: C
%% memo:
\memoanswer{
日光灯或白炽灯发出的都是复色光，其中各种色光的波长不同，产生的衍射条纹间距不同，各单色光相互叠加，成为彩色直条纹，故 C 正确。
}

\item
%% number: 3
%% paperName: 1995年全国高考第 3 题 3 分
%% typeId: 1
%% type: 单选题
%% chapter: 气体性质
%% point: 阿伏伽德罗常数
%% method: 无
%% score: 3
%% degree: 650
%% duplicateId: 0
%% body:
已知铜的密度为 $8.9\times10^{3}\Ukgmc$，原子量为 64。通过估算可知铜中每个铜原子所占的体积为\xzanswer{B}
\fourchoices[answer=B]
{$7\times10^{-6}\Umc$}
{$1\times10^{-29}\Umc$}
{$1\times10^{-26}\Umc$}
{$8\times10^{-24}\Umc$}

%% answer: B
%% memo:
\memoanswer{
铜的摩尔体积 $V=\frac{M}{\rho}$，认为每个原子一个挨着一个排列，每个铜原子所占的体积约为 $V_{0}=\frac{V}{N_{A}}=\frac{M}{\rho N_{A}}=\frac{6.4\times10^{-2}}{8.9\times10^{3}\times6.0\times10^{23}}\Umc\approx1\times10^{-29}\Umc$，故 B 正确。
}

\item
%% number: 4
%% paperName: 1995年全国高考第 4 题 3 分
%% typeId: 1
%% type: 单选题
%% chapter: 原子和原子核
%% point: 半衰期
%% method: 无
%% score: 3
%% degree: 650
%% duplicateId: 0
%% body:
放射性同位素 $\ce{^{24}_{11}Na}$ 的样品经过 $6\Uh$ 后还剩 $\frac{1}{8}$ 没有衰变，它的半衰期是\xzanswer{A}
\fourchoices[answer=A]
{$2\Uh$}
{$1.5\Uh$}
{$1.17\Uh$}
{$0.75\Uh$}

%% answer: A
%% memo:
\memoanswer{
由半衰期公式 $N = N_{0}\left(\frac{1}{2}\right)^{n}$ 知，$\frac{1}{8} = \left(\frac{1}{2}\right)^{3}$，即所经时间 $6\Uh$ 为 3 个半衰期，所以它的半衰期是 $2\Uh$，故 A 正确。
}

\item
%% number: 5
%% paperName: 1995年全国高考第 5 题 3 分
%% typeId: 1
%% type: 单选题
%% chapter: 交流电
%% point: LC振荡回路
%% method: 无
%% score: 3
%% degree: 650
%% duplicateId: 0
%% body:
在 LC 振荡电路中，用以下的哪种办法可以使振荡频率增大一倍？\xzanswer{D}
\fourchoices[answer=D]
{自感 $L$ 和电容 $C$ 都增大一倍}
{自感 $L$ 增大一倍，电容 $C$ 减小一半}
{自感 $L$ 减小一半，电容 $C$ 增大一倍}
{自感 $L$ 和电容 $C$ 都减小一半}

%% answer: D
%% memo:
\memoanswer{
LC 振荡电路的振荡频率为 $f=\frac{1}{2\pi\sqrt{LC}}$，A 项：振荡频率为 $\frac{1}{2}$；B 项：振荡频率不变；C 项：振荡频率不变；D 项：振荡频率增大一倍，故 D 正确。
}

\item
%% number: 6
%% paperName: 1995年全国高考第 6 题 3 分
%% typeId: 1
%% type: 单选题
%% chapter: 光学
%% point: 光电效应
%% method: 无
%% score: 3
%% degree: 850
%% duplicateId: 0
%% body:
在演示光电效应的实验中，原来不带电的一块锌板与灵敏验电器相连。用弧光灯照射锌板时，验电器的指针就张开一个角度，如图 \ref{1995全国06} 所示。这时\xzanswer{B}
\onepicture[num=6, showlabel=false, label=1995全国06, w=7cm]{\includesvg{figs/06}}
\fourchoices[answer=B]
{锌板带正电，指针带负电}
{锌板带正电，指针带正电}
{锌板带负电，指针带正电}
{锌板带负电，指针带负电}

%% answer: B
%% memo:
\memoanswer{
锌板受弧光灯照射，吸收紫外线的光子，发生光电效应，放射出光电子，锌板初始显电中性，放出光电子后显正电性，带正电，验电器与锌板连在一起，也带正电，故 B 正确。
}

\item
%% number: 7
%% paperName: 1995年全国高考第 7 题 3 分
%% typeId: 1
%% type: 单选题
%% chapter: 电路
%% point: 分压与分流
%% method: 无
%% score: 3
%% degree: 650
%% duplicateId: 0
%% body:
两个定值电阻 $R_{1}$、$R_{2}$ 串联后接在输出电压 $U$ 稳定于 $12\UV$ 的直流电源上。有人把一个内阻不是远大于 $R_{1}$、$R_{2}$ 的电压表接在 $R_{1}$ 两端（如图 \ref{1995全国07}），电压表的示数为 $8\UV$。如果把此电压表改接在 $R_{2}$ 的两端，则电压表的示数将\xzanswer{A}
\onepicture[num=7, showlabel=false, label=1995全国07, w=5cm]{figs/07.png}
\fourchoices[answer=A]
{小于 $4\UV$}
{等于 $4\UV$}
{大于 $4\UV$，小于 $8\UV$}
{等于或大于 $8\UV$}

%% answer: A
%% memo:
\memoanswer{
当电压表接在 $R_{1}$ 两端时，其示数为 $8\UV$，则此时电阻 $R_{2}$ 两端的电压为 $4\UV$，将 $R_{1}$ 与 $R_{V}$ 并联后的电阻用 $R_{1V}$ 表示，则 $R_{1V}:R_{2}=8:4=2:1$，即 $R_{1V}=2R_{2}$，由于 $R_{1}>R_{1V}$，则 $R_{1}>2R_{2}$；当电压表接在 $R_{2}$ 两端时，将 $R_{2}$ 与 $R_{V}$ 并联后的电阻用 $R_{2V}$ 表示，则 $R_{2}>R_{2V}$，此时电阻 $R_{1}$ 两端的电压 $U_{1}$ 与电压表示数 $U_{2V}$ 之比 $U_{1}:U_{2V}=R_{1}:R_{2V}>2R_{2}:R_{2V}>2R_{2}:R_{2}=2$，所以电压表的示数将小于 $4\UV$，故 A 正确。
}

\item
%% number: 8
%% paperName: 1995年全国高考第 8 题 3 分
%% typeId: 1
%% type: 单选题
%% chapter: 曲线运动
%% point: 人造地球卫星
%% method: 无
%% score: 3
%% degree: 650
%% duplicateId: 0
%% body:
两颗人造卫星 A、B 绕地球作圆周运动，周期之比为 $T_{A}:T_{B}=1:8$，则轨道半径之比和运动速率之比分别为\xzanswer{D}
\fourchoices[answer=D]
{$R_{A}:R_{B}=4:1$，$v_{A}:v_{B}=1:2$}
{$R_{A}:R_{B}=4:1$，$v_{A}:v_{B}=2:1$}
{$R_{A}:R_{B}=1:4$，$v_{A}:v_{B}=1:2$}
{$R_{A}:R_{B}=1:4$，$v_{A}:v_{B}=2:1$}

%% answer: D
%% memo:
\memoanswer{
地球对卫星的万有引力提供向心力，由万有引力定律和牛顿第二定律得 $G\frac{Mm}{R^{2}}=m\frac{4\pi^{2}}{T^{2}}R$，有 $T^{2}\propto R^{3}$，$R\propto\sqrt[3]{T^{2}}$，由 $T_{A}:T_{B}=1:8$ 得 $R_{A}:R_{B}=1:4$，又 $v=\frac{2\pi R}{T}$，所以 $\frac{v_{A}}{v_{B}}=\frac{R_{A}}{R_{B}}\cdot\frac{T_{B}}{T_{A}}=\frac{1}{4}\cdot\frac{8}{1}=\frac{2}{1}$，故 D 正确。
}

\item
%% number: 9
%% paperName: 1995年全国高考第 9 题 3 分
%% typeId: 1
%% type: 单选题
%% chapter: 运动和力
%% point: 连接体
%% method: 无
%% score: 3
%% degree: 650
%% duplicateId: 0
%% body:
如图 \ref{1995全国09} 所示，质量为 $m$ 的物体 A 放置在质量为 $M$ 的物体 B 上，B 与弹簧相连，它们一起在光滑水平面上作简谐振动，振动过程中 A、B 之间无相对运动。设弹簧的倔强系数为 $k$。当物体离开平衡位置的位移为 $x$ 时，A、B 间摩擦力的大小等于\xzanswer{D}
\onepicture[num=9, showlabel=false, label=1995全国09, w=6cm]{figs/09.png}
\fourchoices[answer=D]
{$0$}
{$kx$}
{$\frac{m}{M}kx$}
{$\frac{m}{M+m}kx$}

%% answer: D
%% memo:
\memoanswer{
因 A、B 之间无相对运动，可把 A、B 看做一个整体，由牛顿第二定律得 $kx=(M+m)a$，解得 $a=\frac{kx}{M+m}$，A 的加速度由 B 对 A 的静摩擦力产生，有 $F_{f}=ma$，即 $F_{f}=\frac{m}{M+m}kx$，故 D 正确。
}

\item
%% number: 10
%% paperName: 1995年全国高考第 10 题 3 分
%% typeId: 1
%% type: 单选题
%% chapter: 电路
%% point: 动态分析
%% method: 无
%% score: 3
%% degree: 450
%% duplicateId: 0
%% body:
在如图 \ref{1995全国10} 所示的电路中，电键 $K_{1}$、$K_{2}$、$K_{3}$、$K_{4}$ 均闭合，C 是极板水平放置的平行板电容器，板间悬浮着一油滴 P。断开哪一个电键后 P 会向下运动？\xzanswer{C}
\onepicture[num=10, showlabel=false, label=1995全国10, w=6cm]{figs/10.png}
\fourchoices[answer=C]
{$K_{1}$}
{$K_{2}$}
{$K_{3}$}
{$K_{4}$}

%% answer: C
%% memo:
\memoanswer{
解法一：分析法，板间悬浮着的油滴 P，受重力和电场力，重力向下，电场力必向上，要使 P 向下运动，必须减小电场力，即减小电场强度，亦即减小平行板电容器的电压，断开 $K_{3}$ 即断开平行板电容器与电源的连接，电容器通过 $R_{3}$ 放电，电压会减小，故 C 正确。

解法二：检验法，检验每一个选项是否符合要求，A 项：断开 $K_{1}$，因所在电路无电流，所以断开 $K_{1}$ 对电路无影响，故 A 错误；B 项：断开 $K_{2}$ 后，电容器的电压从路端电压增大为电源电动势，所以电场力增大，油滴 P 会向上运动，故 B 错误；C 项：断开 $K_{3}$，即断开平行板电容器与电源的连接，电容器通过 $R_{3}$ 放电，电压会减小，电场力减小，油滴 P 会向下运动，故 C 正确；D 项：断开 $K_{4}$，电容器上电量不变，电压不变，电场力不变，油滴 P 仍会悬浮，故 D 错误。
}

\item
%% number: 11
%% paperName: 1995年全国高考第 11 题 3 分
%% typeId: 1
%% type: 单选题
%% chapter: 交流电
%% point: 交流电
%% method: 无
%% score: 3
%% degree: 450
%% duplicateId: 0
%% body:
如图 \ref{1995全国11} 表示一交流电的电流随时间而变化的图像。此交流电流的有效值是\xzanswer{B}
\onepicture[num=11, showlabel=false, label=1995全国11, w=8cm]{figs/11.png}
\fourchoices[answer=B]
{$5\sqrt{2}\UA$}
{$5\UA$}
{$3.5\sqrt{2}\UA$}
{$3.5\UA$}

%% answer: B
%% memo:
\memoanswer{
交流电的有效值的意义：在一个周期的时间内，交流电电流通过电阻 $R$ 产生的热量与恒定电流 $I$ 通过同一电阻 $R$ 产生的热量相等，$I$ 称为 $i$ 的有效值。在一个周期 $T$ 内（$T=0.02\Us$），此交流电在电阻 $R$ 产生的热量为 $Q=I_{1}^{2}\cdot R\cdot\frac{T}{2}+I_{2}^{2}\cdot R\cdot\frac{T}{2}=(4\sqrt{2})^{2}\times0.01\times R+(3\sqrt{2})^{2}\times0.01\times R=0.50R$，又 $Q=I^{2}RT$，所以 $I=5\UA$，故 B 正确。
}

\item
%% number: 12
%% paperName: 1995年全国高考第 12 题 5 分
%% typeId: 2
%% type: 多选题
%% chapter: 运动和力
%% point: 力学单位制
%% method: 无
%% score: 5
%% degree: 850
%% duplicateId: 0
%% body:
在下面列举的物理量单位中，哪些是国际单位制的基本单位？\xzanswer{ABC}
\fourchoices[answer=ABC]
{千克（$\Ukg$）}
{米（$\Um$）}
{开尔文（$\UK$）}
{牛顿（$\UN$）}

%% answer: ABC
%% memo:
\memoanswer{
“牛顿（N）”不是基本单位，是由牛顿第二定律推导的单位，其余是国际单位制基本单位，故 ABC 正确，D 错误。
}

\item
%% number: 13
%% paperName: 1995年全国高考第 13 题 5 分
%% typeId: 1
%% type: 单选题
%% chapter: 机械波
%% point: 波长频率波速
%% method: 无
%% score: 5
%% degree: 750
%% duplicateId: 0
%% body:
关于机械波的概念，下列说法中正确的是\xzanswer{BD}
\fourchoices[answer=BD]
{质点振动的方向总是垂直于波传播的方向}
{简谐波沿长绳传播，绳上相距半个波长的两质点振动位移的大小相等}
{任一振动质点每经过一个周期沿波的传播方向移动一个波长}
{相隔一个周期的两时刻，简谐波的图像相同}

%% answer: BD
%% memo:
\memoanswer{
由机械波的概念可知，由于机械波按质点振动方向与波传播方向间的关系分为横波和纵波，质点振动方向与波的传播方向垂直的波为横波，质点振动方向与波的传播方向在一条直线上的波是纵波，故 A 错误；简谐波在某时刻的波动图象如图所示，相距半个波长的两质点如图中 $P_{1}$ 和 $P_{2}$，在同一时刻会振动到距平衡位置相等距离的位置，即位移大小相等，方向相反，故 B 正确；振动质点只在其平衡位置附近做往返运动，而不随波移动，故 C 错误；由于每个质点都在做简谐运动，因此每个质点经过一个周期后均会重复原来的运动状态，所以相隔一个周期的两时刻，简谐波的图象相同，故 D 正确。

\onepicture[w=8cm]{figs/13_an.jpg}
}

\item
%% number: 14
%% paperName: 1995年全国高考第 14 题 5 分
%% typeId: 2
%% type: 多选题
%% chapter: 电场
%% point: 等势面
%% method: 无
%% score: 5
%% degree: 650
%% duplicateId: 0
%% body:
在静电场中\xzanswer{CD}
\fourchoices[answer=CD]
{电场强度处处为零的区域内，电势也一定处处为零}
{电场强度处处相同的区域内，电势也一定处处相同}
{电场强度的方向总是跟等势面垂直的}
{沿着电场强度的方向，电势总是不断降低的}

%% answer: CD
%% memo:
\memoanswer{
电势具有相对性，零电势的选取是任意的，所以电场强度为零的区域电势不一定为零，故 A 错误；在匀强电场中，电场强度处处相同，但电势不相同，故 B 错误；由电场的性质知 CD 正确。
}

\item
%% number: 15
%% paperName: 1995年全国高考第 15 题 5 分
%% typeId: 2
%% type: 多选题
%% chapter: 光学
%% point: 光的折射
%% method: 无
%% score: 5
%% degree: 650
%% duplicateId: 0
%% body:
已知介质对某单色光的临界角为 $\theta$，则\xzanswer{ABC}
\fourchoices[answer=ABC]
{该介质对此单色光的折射率等于 $\frac{1}{\sin\theta}$}
{此单色光在该介质中的传播速度等于 $c\sin\theta$（$c$ 是真空中的光速）}
{此单色光在该介质中的波长是在真空中波长的 $\sin\theta$ 倍}
{此单色光在该介质中的频率是在真空中频率的 $\frac{1}{\sin\theta}$ 倍}

%% answer: ABC
%% memo:
\memoanswer{
由临界角的概念知 $\sin\theta=\frac{1}{n}$，所以折射率为 $n=\frac{1}{\sin\theta}$，故 A 正确；由折射率与光速的关系 $n=\frac{c}{v}$ 知，$v=\frac{c}{n}=c\sin\theta$，故 B 正确；由 $n=\frac{c}{v}=\frac{\lambda_{0}}{\lambda}$ 知 $\frac{\lambda}{\lambda_{0}}=\frac{1}{n}=\sin\theta$，故 C 正确；光从一种介质到另一种介质频率不变，故 D 错误。
}

\item
%% number: 16
%% paperName: 1995年全国高考第 16 题 5 分
%% typeId: 2
%% type: 多选题
%% chapter: 磁场
%% point: 带电粒子在磁场中的运动
%% method: 无
%% score: 5
%% degree: 650
%% duplicateId: 0
%% body:
两个粒子，带电量相等，在同一匀强磁场中只受磁场力而作匀速圆周运动。\xzanswer{BC}
\fourchoices[answer=BC]
{若速率相等，则半径必相等}
{若质量相等，则周期必相等}
{若动量大小相等，则半径必相等}
{若动能相等，则周期必相等}

%% answer: BC
%% memo:
\memoanswer{
带电粒子在匀强磁场中做匀速圆周运动的半径 $R=\frac{mv}{qB}$，周期 $T=\frac{2\pi m}{qB}$。由于 $q$、$B$ 相同，故 $m$ 相等时 $T$ 相等，$mv$ 大小相等时 $R$ 相等，故 BC 正确。
}

\item
%% number: 17
%% paperName: 1995年全国高考第 17 题 5 分
%% typeId: 2
%% type: 多选题
%% chapter: 功和能
%% point: 功能原理
%% method: 无
%% score: 5
%% degree: 450
%% duplicateId: 0
%% body:
一粒钢珠从静止状态开始自由下落，然后陷入泥潭中。若把在空中下落的过程称为过程 Ⅰ，进入泥潭直到停住的过程称为过程 Ⅱ，则\xzanswer{AC}
\fourchoices[answer=AC]
{过程 Ⅰ 中钢珠动量的改变量等于重力的冲量}
{过程 Ⅱ 中阻力的冲量的大小等于过程 Ⅰ 中重力冲量的大小}
{过程 Ⅱ 中钢珠克服阻力所做的功等于过程 Ⅰ 与过程 Ⅱ 中钢珠所减少的重力势能之和}
{过程 Ⅱ 中损失的机械能等于过程 Ⅰ 中钢珠所增加的动能}

%% answer: AC
%% memo:
\memoanswer{
在过程Ⅰ中钢球受到的合外力就是重力，由动量定理知在过程Ⅰ中小球动量的变化等于重力的冲量，故 A 正确；对于Ⅰ和Ⅱ全过程运用动量定理得 $I_{1} + I_{2} - I_{3} = 0$，其中 $I_{1}$、$I_{2}$ 分别表示在过程Ⅰ、Ⅱ中重力的冲量，$I_{3}$ 表示在过程Ⅱ中阻力的冲量，得 $I_{3} = I_{1} + I_{2}$，故 B 错误；在过程Ⅰ和Ⅱ全过程应用功能关系可知，在过程Ⅰ和Ⅱ中钢球减少的重力势能全部用来克服阻力做功转化为内能，故 C 正确，D 错误。
}

\item
%% number: 18
%% paperName: 1995年全国高考第 18 题 5 分
%% typeId: 1
%% type: 单选题
%% chapter: 机械振动
%% point: 弹簧振子
%% method: 无
%% score: 5
%% degree: 450
%% duplicateId: 0
%% body:
一弹簧振子作简谐振动，周期为 $T$，则\xzanswer{C}
\fourchoices[answer=C]
{若 $t$ 时刻和（$t+\Delta t$）时刻振子运动位移的大小相等、方向相同，则 $\Delta t$ 一定等于 $T$ 的整数倍}
{若 $t$ 时刻和（$t+\Delta t$）时刻振子运动速度的大小相等、方向相反，则 $\Delta t$ 一定等于 $T/2$ 的整数倍}
{若 $\Delta t=T$，则在 $t$ 时刻和（$t+\Delta t$）时刻振子运动的加速度一定相等}
{若 $\Delta t=T/2$，则在 $t$ 时刻和（$t+\Delta t$）时刻弹簧的长度一定相等}

%% answer: C
%% memo:
\memoanswer{
如图所示为做简谐运动的物体的振动图象，由图可知，在同一个全振动，即 1 个周期内存在两个时刻 $t_{1}$、$t_{2}$，此时质点在平衡位置同侧的同一位置，即位移大小相等、方向相同，但 $\Delta t = t_{2} - t_{1}$ 不等于 $T$ 的整数倍，故 A 错误；由图可知，当 $t_{2} - t_{1} < T/2$ 时，也会存在质点的运动速度大小相等，但一时刻为向远离平衡位置的方向振动，另一时刻为向靠近平衡位置的方向振动，即方向相反，故 B 错误；由图可知，当两时刻 $t_{1}$、$t_{2}$ 相差一个周期时，会有质点振动情况完全相同，因此加速度也应相等，故 C 正确；$t_{1}$、$t_{2}$ 两时刻相差 $T/2$ 时，由于质点位于平衡位置两侧，即一时刻是压缩弹簧，另一时刻是拉伸弹簧，因此弹簧的长度不相等，故 D 错误。

\onepicture[w=9cm]{figs/18_an.jpg}
}

\item
%% number: 19
%% paperName: 1995年全国高考第 19 题 5 分
%% typeId: 3
%% type: 填空题
%% chapter: 功和能
%% point: 动能定理
%% method: 无
%% score: 5
%% degree: 750
%% duplicateId: 0
%% body:
一人坐在雪橇上，从静止开始沿着高度为 $15\Um$ 的斜坡滑下，到达底部时速度为 $10\Ums$。人和雪橇的总质量为 $60\Ukg$，下滑过程中克服阻力做的功等于\tkanswer[2.5cm]{$6000\UJ$}（取 $g=10\Umsq$）。

%% answer: $6000\UJ$
\jdanswer{
$6000\UJ$
}

%% memo:
\memoanswer{
对整个运动过程，由动能定理得 $mgh - W_{f} = \frac{1}{2}mv^{2} - 0$，解得 $W_{f} = 6000\UJ$。
}

\item
%% number: 20
%% paperName: 1995年全国高考第 20 题 5 分
%% typeId: 3
%% type: 填空题
%% chapter: 原子和原子核
%% point: 玻尔模型
%% method: 无
%% score: 5
%% degree: 750
%% duplicateId: 0
%% body:
如图 \ref{1995全国20} 给出氢原子最低的四个能级。氢原子在这些能级之间跃迁所辐射的光子的频率最多有\tkanswer[1cm]{6}种，其中最小的频率等于\tkanswer[2.5cm]{$1.6\times10^{14}\UHz$}。
\onepicture[num=20, showlabel=false, label=1995全国20, align=right, w=5cm]{\includesvg{figs/20}}

%% answer: 6，$1.6\times10^{14}\UHz$
\jdanswer{
6，$1.6\times10^{14}\UHz$
}

%% memo:
\memoanswer{
氢原子最低的四个能级之间的辐射跃迁，如图所示，有 6 种可能，对应发射光子的能量有 6 种，对应的频率 $\nu$ 也有 6 种，其中从 $n=4$ 跃迁到 $n=3$ 能级时原子辐射能量最小，即光子频率最小，由玻尔理论得 $\nu=\frac{E_{4}-E_{3}}{h}=1.6\times10^{14}\UHz$。

\onepicture[w=6cm]{figs/20_an.jpg}
}

\item
%% number: 21
%% paperName: 1995年全国高考第 21 题 5 分
%% typeId: 3
%% type: 填空题
%% chapter: 运动和力
%% point: 牛顿第二定律
%% method: 无
%% score: 5
%% degree: 650
%% duplicateId: 0
%% body:
已知质量为 $m$ 的木块在大小为 $T$ 的水平拉力作用下沿粗糙水平地面作匀加速直线运动，加速度为 $a$，则木块与地面之间的滑动摩擦因数为\tkanswer[2.5cm]{$\frac{T-ma}{mg}$}。若在木块上再施加一个与水平拉力 $T$ 在同一竖直平面内的推力，而不改变木块速度的大小和方向，则此推力与水平拉力 $T$ 的夹角为\tkanswer[3cm]{$\arctan\frac{mg}{T-ma}$}。

%% answer: $\frac{T-ma}{mg}$，$\arctan\frac{mg}{T-ma}$
\jdanswer{
$\frac{T-ma}{mg}$，$\arctan\frac{mg}{T-ma}$
}

%% memo:
\memoanswer{
木块受拉力 $T$ 作用做匀加速直线运动，设此时木块受的滑动摩擦力是 $F_{f1}$，地面对木块的支持力为 $N_{1}$，由牛顿第二定律得 $T - F_{f1} = ma$，又 $N_{1} = mg$，$F_{f1} = \mu N_{1}$，联立解得 $\mu = \frac{T - ma}{mg}$。

设再施加的推力为 $F$，与水平拉力 $T$ 的夹角为 $\alpha$，此时木块受地面的支持力为 $N_{2}$，摩擦力为 $F_{f2}$，将 $F$ 分解为水平分力 $F_{1}$ 及竖直分力 $F_{2}$，则 $F_{1}=F\cos\alpha$，$F_{2}=F\sin\alpha$，由牛顿第二定律得 $T+F_{1}-F_{f2}=ma$，又 $T-F_{f1}=ma$，联立得 $F_{1}-F_{f2}=-F_{f1}$，又 $F_{f2}=\mu N_{2}$，$N_{2}=F_{2}+mg$，$F_{f1}=\mu mg$，代入得 $F_{1}=\mu F_{2}$，解得 $\tan\alpha=\frac{1}{\mu}$，故 $\alpha=\arctan\frac{1}{\mu}=\arctan\frac{mg}{T-ma}$。
}

\item
%% number: 22
%% paperName: 1995年全国高考第 22 题 5 分
%% typeId: 3
%% type: 填空题
%% chapter: 电场
%% point: 电势能电势差电场力做功
%% method: 无
%% score: 5
%% degree: 450
%% duplicateId: 0
%% body:
a、b 和 c 表示点电荷的电场中的三个等势面，它们的电势分别为 $\varphi$、$\frac{2}{3}\varphi$ 和 $\frac{1}{4}\varphi$。一带电粒子从等势面 a 上某处由静止释放后，仅受电场力作用而运动。已知它经过等势面 b 时的速率为 $v$，则它经过等势面 c 时的速率为\tkanswer[2cm]{$1.5v$}。
\onepicture[num=22, showlabel=false, label=1995全国22, align=right, w=4.5cm]{\includesvg{figs/22}}

%% answer: $1.5v$
\jdanswer{
$1.5v$
}

%% memo:
\memoanswer{
带电粒子从等势面 a 运动到等势面 b 的过程中，由动能定理得 $q\left(U-\frac{2}{3}U\right)=\frac{1}{2}mv^{2}$，同样粒子从 a 到 c 的过程中，设粒子经过等势面 c 时的速率为 $v_{c}$，有 $q\left(U-\frac{1}{4}U\right)=\frac{1}{2}mv_{c}^{2}$，解得 $v_{c}=1.5v$。
}

\item
%% number: 23
%% paperName: 1995年全国高考第 23 题 5 分
%% typeId: 7
%% type: 作图题
%% chapter: 光学
%% point: 透镜
%% method: 无
%% score: 5
%% degree: 550
%% duplicateId: 0
%% body:
图中 AB 表示一直立的平面镜，$P_{1}P_{2}$ 是水平放置的米尺（有刻度的一面朝着平面镜），MN 是屏，三者互相平行。屏 MN 上的 ab 表示一条竖直的缝（即 a、b 之间是透光的）。某人眼睛紧贴米尺上的小孔 S（其位置见图），可通过平面镜看到米尺的一部分刻度。试在本题的图上用三角板作图求出可看到的部位，并在 $P_{1}P_{2}$ 上把这部分涂以标志。
\drawpicanswer{\onepicture[align=right, w=6cm]{figs/23.png}}{\onepicture[align=right, w=10cm]{figs/23_an.jpg}}

%% answer: 如图
\jdanswer{
如图
}

%% memo:
\memoanswer{
由平面镜成像规律，确定眼睛像（小孔 S）的位置 $S'$，作直线 $Sa$ 并延长找到与平面镜的交点，再作直线 $S'a$ 并延长，找到与 $P_{1}P_{2}$ 的交点 $a'$；作直线 $S'b$ 并延长，找到与 $P_{1}P_{2}$ 的交点 $b'$，再由光的反射定律补画出相应的光路，故可看到的部分为 $a'b'$。

另解：由平面镜成像的规律分别作出 $P_{1}P_{2}$ 和 MN 在 AB 中的虚像如图所示，由图可见，作直线 $Sa$ 并延长找到与 $P_{1}P_{2}'$ 的交点 $a'$，再作直线 $Sb'$ 并延长找到与 $P_{1}P_{2}'$ 的交点 $b'$，再由光的反射定律和平面镜成像规律补画出相应的光路，故可看到的部分为 $a'b'$。
}

\item
%% number: 24
%% paperName: 1995年全国高考第 24 题 6 分
%% typeId: 4
%% type: 实验题
%% chapter: 电磁感应
%% point: 验证电磁感应实验
%% method: 无
%% score: 6
%% degree: 650
%% duplicateId: 0
%% body:
在研究电磁感应现象的实验中所用的器材如图所示，它们是：
（1）电流计，（2）直流电源，（3）带铁心的线圈 A，（4）线圈 B，（5）电键，（6）滑动变阻器。（用来控制电流以改变磁场强弱）

试按实验的要求在实物图上连线。（图中已连好一根导线）

若连接滑动变阻器的两根导线接在线柱 C 和 D 上，而在电键刚闭合时电流计指针右偏，则电键闭合后滑动变阻器的滑动触头向接线柱 C 移动时，电流计指针将\tkanswer[1.5cm]{左偏}。（填“左偏”、“右偏”或“不偏”）
\drawpicanswer{\onepicture[align=right, w=8cm]{figs/24.png}}{\onepicture[align=right, w=8cm]{figs/24_an.jpg}}

%% answer: 如图所示；左偏
\jdanswer{
如图所示；左偏。
}

%% memo:
\memoanswer{
如图所示为一种连接方法，只要将 B 线圈与电流表串联起来，将滑动变阻器、电池、开关和 A 线圈串联起来即可。当电键闭合时，电流增大，磁场增强，电流表指针右偏；当滑动片向右滑时，滑动变阻器接在电路中的电阻值增大，电路中电流减小，磁场减弱，B 线圈中感应电流的方向应与前面情形相反，故电流表指针向左偏。
}

\item
%% number: 25
%% paperName: 1995年全国高考第 25 题 6 分
%% typeId: 3
%% type: 填空题
%% chapter: 曲线运动
%% point: 平抛运动
%% method: 无
%% score: 6
%% degree: 450
%% duplicateId: 0
%% body:
在研究平抛物体运动的实验中，用一张印有小方格的纸记录轨迹，小方格的边长 $l=1.25\Ucm$。若小球在平抛运动途中的几个位置如图 \ref{1995全国25} 中的 a、b、c、d 所示，则小球平抛的初速度的计算式为 $v_{0}=$\tkanswer[2cm]{$2\sqrt{gl}$}（用 $l$、$g$ 表示），其值是\tkanswer[2.5cm]{$0.70\Ums$}。（取 $g=9.8\Umsq$）
\onepicture[num=25, showlabel=false, label=1995全国25, align=right, w=5cm]{\includesvg{figs/25}}

%% answer: $2\sqrt{gl}$，$0.70\Ums$
\jdanswer{
$2\sqrt{gl}$，$0.70\Ums$
}

%% memo:
\memoanswer{
平抛运动水平方向的分运动是匀速直线运动，竖直方向的分运动是自由落体运动。从图中看出，a、b、c、d 点间的水平间距均相等，间距 $x = 2l$，因此这 4 个点是等时间间隔点，竖直方向两段相邻位移之差 $\Delta y = l$，又 $\Delta y = gT^{2}$，$v_{0} = \frac{2l}{T}$，解出 $v_{0} = 2\sqrt{lg}$，代入 $l = 1.25\Ucm$，$g = 9.8\Umsq$，得 $v_{0} = 0.70\Ums$。
}

\item
%% number: 26
%% paperName: 1995年全国高考第 26 题 5 分
%% typeId: 4
%% type: 实验题
%% chapter: 电路
%% point: 多用表的使用
%% method: 无
%% score: 5
%% degree: 750
%% duplicateId: 0
%% body:
某人用万用电表按正确步骤测量一电阻阻值，指针指示位置如图 \ref{1995全国26} 所示，则这电阻值是\tkanswer[2.5cm]{$1.2\times10^{3}\UO$}。如果要用这万用电表测量一个约 $200\UO$ 的电阻，为了使测量比较精确，选择开关应选的欧姆挡是\tkanswer[2cm]{$\times10$}。
\onepicture[num=26, showlabel=false, label=1995全国26, align=right, w=7cm]{figs/26.png}

%% answer: $1.2\times10^{3}\UO$；$\times10$
\jdanswer{
$1.2\times10^{3}\UO$；$\times10$
}

%% memo:
\memoanswer{
表针示数是 $12\UO$，挡位是 $\times100\UO$，故测得此电阻值为 $12 \times 100\UO = 1.2 \times 10^{3}\UO$。若测约 $200\UO$ 的电阻，中值电阻在 $15 \sim 20\UO$ 之间，取 $20\UO$，应选 $\times10$ 挡。
}

\item
%% number: 27
%% paperName: 1995年全国高考第 27 题 6 分
%% typeId: 6
%% type: 计算题
%% chapter: 光学
%% point: 透镜
%% method: 无
%% score: 6
%% degree: 750
%% duplicateId: 0
%% body:
一发光点 S 位于焦距为 $12\Ucm$ 的薄凸透镜的主轴上。当 S 沿垂直于主轴的方向移动 $1.5\Ucm$ 时，它的像点 $S'$ 移动 $0.5\Ucm$。求移动前发光点 S 到像点 $S'$ 的距离。

%% answer: $l = 64\Ucm$
\jdanswer{
$l = 64\Ucm$
}

%% memo:
\memoanswer{
解：用 $h$ 和 $h'$ 分别表示 $S$ 和 $S'$ 移动的距离，用 $l$ 表示 $S$ 和 $S'$ 未移动时的距离，则有

$\frac{v}{u} = \frac{h'}{h}$ ①，

$l = u + v$ ②

根据透镜成像公式

$\frac{1}{u} + \frac{1}{v} = \frac{1}{f}$ ③，

由①②③式并代入数据可解得

$l = 64\Ucm$

评分标准：全题 6 分，写出①式给 2 分，写出②式给 2 分，写出③式给 1 分。得出正确结果再给 1 分。
}

\item
%% number: 28
%% paperName: 1995年全国高考第 28 题 10 分
%% typeId: 6
%% type: 计算题
%% chapter: 电磁感应
%% point: 电磁感应电路分析
%% method: 无
%% score: 10
%% degree: 550
%% duplicateId: 0
%% body:
两根相距 $d=0.20\Um$ 的平行金属长导轨固定在同一水平面内，并处于竖直方向的匀强磁场中，磁场的磁感应强度 $B=0.2\UT$，导轨上面横放着两条金属细杆，构成矩形回路，每条金属细杆的电阻为 $r=0.25\UO$，回路中其余部分的电阻可不计。已知两金属细杆在平行于导轨的拉力的作用下沿导轨朝相反方向匀速平移，速度大小都是 $v=5.0\Ums$，如图 \ref{1995全国28} 所示。不计导轨上的摩擦。
\begin{enumerate}
	\item
	求作用于每条金属细杆的拉力的大小。
	\item
	求两金属细杆在间距增加 $0.40\Um$ 的滑动过程中共产生的热量。
\end{enumerate}
\onepicture[num=28, showlabel=false, label=1995全国28, align=right, w=7cm]{figs/28.png}

%% answer: （1）$F_{1}=F_{2}=3.2\times10^{-2}\UN$；（2）$Q=1.28\times10^{-2}\UJ$
\jdanswer{
\begin{enumerate}
	\item
	$F_{1}=F_{2}=3.2\times10^{-2}\UN$

	\item
	$Q=1.28\times10^{-2}\UJ$
\end{enumerate}
}

%% memo:
\memoanswer{
解：（1）当两金属杆都以速度 $v$ 匀速滑动时，每条金属杆中产生的感应电动势分别为 $E_{1}=E_{2}=Bdv$ ①

由闭合电路的欧姆定律，回路中的电流强度

$I=\frac{E_{1}+E_{2}}{2r}$ ②，

因拉力与安培力平衡，作用于每根金属杆的拉力的大小为

$F_{1}=F_{2}=BId$ ③

由①②③式并代入数据得

$F_{1}=F_{2}=\frac{B^{2}d^{2}v}{r}=3.2\times10^{-2}\UN$

（2）设两金属杆之间增加的距离为 $\Delta L$，则两金属杆共产生的热量

$Q = I^{2} \cdot 2r \cdot \frac{\Delta L}{2v}$ ④，

代入数据得 $Q = 1.28 \times 10^{-2}\UJ$

评分标准：全题 10 分。第一问 6 分：求出①式给 1 分，求出②③式各得 2 分，结果正确再给 1 分。第二问 4 分：求出④式给 3 分，结果正确再给 1 分。若用 $Q = F_{1} \Delta L$ 代替④式也同样给分。
}

\item
%% number: 29
%% paperName: 1995年全国高考第 29 题 12 分
%% typeId: 6
%% type: 计算题
%% chapter: 气体性质
%% point: 气体多过程
%% method: 无
%% score: 12
%% degree: 450
%% duplicateId: 0
%% body:
一个质量可不计的活塞将一定量的理想气体封闭在上端开口的直立圆筒形气缸内，活塞上堆放着铁砂，如图 \ref{1995全国29} 所示。最初活塞搁置在气缸内壁的固定卡环上，气体柱的高度为 $H_{0}$，压强等于大气压强 $p_{0}$。现对气体缓慢加热，当气体温度升高了 $\Delta T=60\UK$ 时，活塞（及铁砂）开始离开卡环而上升。继续加热直到气柱高度为 $H_{1}=1.5H_{0}$。此后，在维持温度不变的条件下逐渐取走铁砂，直到铁砂全部取走时，气柱高度变为 $H_{2}=1.8H_{0}$，求此时气体的温度。（不计活塞与气缸之间的摩擦）
\onepicture[num=29, showlabel=false, label=1995全国29, align=right, w=5cm]{\includesvg{figs/29}}

%% answer: $T_{2}=540\UK$
\jdanswer{
$T_{2}=540\UK$
}

%% memo:
\memoanswer{
第一种解法：设气体最初温度为 $T_{0}$，则活塞刚离开卡环时温度为 $T_{0} + \Delta T$，压强 $p_{1}$。由等容升温过程得

$\frac{T_{0}+\Delta T}{T_{0}}=\frac{p_{1}}{p_{0}}$ ①，

设气柱高度为 $H_{1}$ 时温度为 $T_{1}$，由等压升温过程得

$\frac{T_{1}}{T_{0}+\Delta T}=\frac{H_{1}}{H_{0}}$ ②，

设气柱高度为 $H_{2}$ 时温度为 $T_{2}$，由等温膨胀过程（$T_{2}=T_{1}$）得

$\frac{p_{0}}{p_{1}}=\frac{H_{1}}{H_{2}}$ ③，

由①和③两式求得

$\frac{T_{0}+\Delta T}{T_{0}}=\frac{H_{2}}{H_{1}}$ ④，解得 $T_{0}=\frac{H_{1}}{H_{2}-H_{1}}\Delta T$ ⑤，

由②和④两式得

$\frac{T_{1}}{T_{0}}=\frac{H_{2}}{H_{0}}$，或 $T_{1}=\frac{H_{2}}{H_{0}}T_{0}$ ⑥，

将⑤式代入⑥式，并利用 $T_{2}=T_{1}$ 得

$T_{2}=T_{1}=\frac{H_{1}H_{2}}{H_{0}(H_{2}-H_{1})}\Delta T$ ⑦，

代入数据得 $T_{2}=540\UK$

评分标准：全题 12 分。求得①、②、③式各给 3 分。正确求得⑦式给 2 分，结果正确再给 1 分（若利用①、②、③式得出正确结果而未写⑦式，也给这 3 分）。

第二种解法：设气体最初温度为 $T_{0}$，则活塞刚离开卡环时温度为 $T_{0} + \Delta T$。设气柱高度为 $H_{1}$ 时温度为 $T_{1}$，高度为 $H_{2}$ 时温度为 $T_{2}$。

由等压升温过程得

$\frac{H_{0}}{T_{0}+\Delta T}=\frac{H_{1}}{T_{1}}$ ①，

联系初态和终态，由理想气体状态方程得

$\frac{H_{0}}{T_{0}}=\frac{H_{2}}{T_{2}}$ ②，

利用 $T_{1}=T_{2}$，由①、②两式解得

$T_{2}=\frac{H_{1}H_{2}}{H_{0}(H_{2}-H_{1})}\Delta T$ ③

代入数值得 $T_{2}=540\UK$

评分标准：全题 12 分。求得①式给 4 分；求得②式给 5 分；正确求得③式给 2 分，结果正确再给 1 分（若利用①、②式得出正确结果而未写③式的，也给这 3 分）。
}

\item
%% number: 30
%% paperName: 1995年全国高考第 30 题 12 分
%% typeId: 6
%% type: 计算题
%% chapter: 运动和力
%% point: 动量守恒定律
%% method: 无
%% score: 12
%% degree: 250
%% duplicateId: 0
%% body:
如图 \ref{1995全国30} 所示，一排人站在沿 $x$ 轴的水平轨道旁，原点 $O$ 两侧的人的序号都记为 $n$（$n=1$，2，3$\ldots$）。每人只有一个沙袋，$x>0$ 一侧的每个沙袋质量为 $m=14\Ukg$，$x<0$ 一侧的每个沙袋质量 $m'=10\Ukg$。一质量为 $M=48\Ukg$ 的小车以某初速度从原点出发向正 $x$ 方向滑行。不计轨道阻力。当车每经过一人身旁时，此人就把沙袋以水平速度 $u$ 朝与车速相反的方向沿车面扔到车上，$u$ 的大小等于扔此袋之前的瞬间车速大小的 $2n$ 倍。（$n$ 是此人的序号数）
\begin{enumerate}
	\item
	空车出发后，车上堆积了几个沙袋时车就反向滑行？
	\item
	车上最终有大小沙袋共多少个？
\end{enumerate}
\onepicture[num=30, showlabel=false, label=1995全国30, align=right, w=9cm]{figs/30.png}

%% answer: （1）车上堆积 3 个沙袋后车就反向滑行；（2）最终共有大小沙袋 11 个
\jdanswer{
\begin{enumerate}
	\item
	车上堆积 3 个沙袋后车就反向滑行。

	\item
	最终共有大小沙袋 11 个。
\end{enumerate}
}

%% memo:
\memoanswer{
解：（1）在小车朝正 $x$ 方向滑行的过程中，第 $(n-1)$ 个沙袋扔到车上后的车速为 $v_{n-1}$，第 $n$ 个沙袋扔到车上后的车速为 $v_{n}$，由动量守恒定律有

$[M+(n-1)m]v_{n-1}-2nmv_{n-1}=(M+nm)v_{n}$

解得 $v_{n}=\frac{M-(n+1)m}{M+nm}v_{n-1}$ ①，

小车反向运动的条件是 $v_{n-1}>0$，$v_{n}<0$，即

$M-nm>0$ ②

$M-(n+1)m\leq0$ ③

代入数据得

$n<\frac{M}{m}=\frac{48}{14}$，$n\geq\frac{M}{m}-1=\frac{34}{14}$

$n$ 应为整数，故 $n=3$，即车上堆积 3 个沙袋后车就反向滑行。

（2）车自反向滑行直到接近 $x < 0$ 一侧第 1 人所在位置时，车速保持不变，而车的质量为 $M + 3m$。若在朝负 $x$ 方向滑行过程中，第 $(n-1)$ 个沙袋扔到车上后车速为 $v_{n-1}'$，第 $n$ 个沙袋扔到车上后车速为 $v_{n}'$，现取在图中向左的方向（负 $x$ 方向）为速度 $v_{n}'$、$v_{n-1}'$ 的正方向，则由动量守恒定律有

$[M+3m+(n-1)m']v_{n-1}' - 2nm'v_{n-1}' = (M+3m+nm')v_{n}'$

解得 $v_{n}'=\frac{M+3m-(n+1)m'}{M+3m+nm'}v_{n-1}'$ ④，

车不再向左滑行的条件是

$v_{n-1}' > 0$，$v_{n}' \leqslant 0$

即 $M + 3m - nm' > 0$ ⑤

$M + 3m - (n + 1)m' \leqslant 0$ ⑥

代入数据得

$n<\frac{M+3m}{m'}=9$，$n\geq\frac{M+3m}{m'}-1=8$，解得 $8\leq n<9$ ⑦

$n=8$ 时，车停止滑行，即在 $x<0$ 一侧第 8 个沙袋扔到车上后车就停住。故车上最终共有大小沙袋 $3+8=11$ 个。

评分标准：全题 12 分。第（1）问 4 分：求得①式给 2 分，正确分析车反向滑行条件并求得反向时车上沙袋数再给 2 分。（若未求得①式，但求得第 1 个沙袋扔到车上后的车速，正确的也给 2 分。通过逐次计算沙袋扔到车上后的车速，并求得车开始反向滑行时车上沙袋数，也再给 2 分。）第（2）问 8 分：求得④式给 3 分，⑤式给 1 分，⑥式给 2 分。求得⑦式给 1 分。得到最后结果再给 1 分。（若未列出⑤、⑥两式，但能正确分析并得到左侧 $n=8$ 的结论，也可给上述⑤、⑥、⑦式对应的 4 分。）
}

\end{enumerate}

\end{document}
